% R152 source proposal; provenance and exact prior text remain in the proposal manifest.
\section{A grounded conditional learner with distinct output interfaces}
The retention question now needs a mechanism: what information would guide
the movement, and what independent observation could show that the guidance
was useful? Follow one action through the proposed measurements and learner
before adding a writer or attributing a gain to biological teaching.

Start with a movement we can measure well. For the proposed snap showcase,
the arm begins extended with the palm down. The elbow and forearm swing and
turn; the wrist and hand turn upward during the snap, then move down to rest.
The model must account for that whole movement, including the preparation,
contact and return. Finger pairing, organism-specific anatomy and timing still
need a real-reference series; a human-looking animation cannot establish that
a particular monkey can perform the same gesture.

The cup reach supplies a simpler precision comparison: relate recorded signals
to measured movement before attempting the coordinated snap. Each task still
needs its own anatomy, contact, timing and independent comparison.

That comparison connects three interacting groups: task-linked brain
observations, named-muscle observations and a musculoskeletal simulation.
A learner would relate the observations to the mechanics; independent movement
and contact measurements would test its predictions. EEG records scalp
electrical activity and surface EMG records muscle electrical activity.
A separately recorded nerve or spinal signal needs its own declared channel.
The proposed apparatus uses two arm sleeves, scalp EEG and synchronized,
calibrated left and right tripod views. Vendor, channel, bandwidth and clock
qualification remain open.

Intent to action is the target. By understanding and measuring the intermediary,
we aim to make that mapping as close to one-to-one as possible within a declared
task. Specify what counts as the intended action, the participant and session,
state and context, permitted history, tolerated error and independent reference.
The test then asks whether observing the middle reduces ambiguity and improves
prediction or qualified control. It does not assume a unique hidden intention
or claim that the best possible mapping has already been achieved.

The learner's first test would compare its predictions with recorded brain
and muscle activity and with measurements of where the hand moved and when.
Sensory feedback and physical action would add a further question: what did
each separate output interface actually deliver? A successful prediction
cannot answer that question, and EEG and EMG recorders alone cannot write
sensation. For this proposal, no paired trial, trained model, delivered sensory
output or physical actuation is established. Sources and permissions, physical
and computational behaviour, clinical interfaces, reference comparisons and
held-out tests still require independent assessment.

The retained LINK comparison belongs to a different experiment. Its 18 runs
and 720 saved epochs are not paired EEG/EMG/movement trials; both model-admission
gates and all five benefit claims failed. The prospective methods below would
have to supply the missing paired evidence and separately test recording,
prediction, simulated movement, sensory delivery and physical action. A result
for one of those jobs would not establish the others.

\subsection{Form, function and action need a measured middle}
A structural map constrains what a system could do. A NeuroConnectome maps
specified neural connections; its organism, coverage, resolution and source
would have to be declared. The signals passing through those connections
still have to be observed or inferred, with the difference made explicit.
To predict the next action, we need the processes that turn structural
constraints into an outcome. The proposed programme follows form to function
to action, while intent to action names the task-level translation we want
to improve.

DNA to RNA and RNA to protein put a process between a starting constraint and
its downstream product. Regulatory context, RNA processing and cell state
change what that process produces; protein structure alone does not determine
the behaviour of a whole cell. AlphaGenome, RNA or cell-state models and
AlphaFold illustrate these different modelling jobs. AlphaTranscriptome names
a desired predictive capability here. Connecting the predictors would still
leave system dynamics and the environment to account for, with experimental
outcomes needed to test the combined prediction.
% Source-role pointers: alphagenome2026, scgpt, alphafold3; bounded owner checks.
\par\noindent \href{https://www.nature.com/articles/s41586-025-10014-0}{AlphaGenome regulatory prediction}~\cite{alphagenomeRetained};
\href{https://www.nature.com/articles/s41592-024-02201-0}{scGPT cell-state modelling}~\cite{scgptRetained};
\href{https://www.nature.com/articles/s41586-024-07487-w}{AlphaFold 3 structure prediction}~\cite{abramson2024}.
These provide modelling examples, not evidence of an integrated cell programmer.
Our account of AlphaGenome and scGPT is limited to their reported modelling roles;
their complete methods have not been independently assessed here.

In the biological pathway, Brain2Spine and Spine2Body name proposed links to
test; their composition is the Brain2Body aim. The spinal cord integrates
descending and sensory signals and contributes local coordination, reflex
regulation and adaptation. This is a different mechanism from RNA processing,
but it serves the same explanatory purpose: account for the context-dependent
process between an initial constraint and its outcome. Supported spinal
observations would test that account. Where they are absent, the internal
states remain inferred or simulated. No programme spinal intervention or qualified spinal-writing interface is demonstrated here.
\par\noindent \href{https://pubmed.ncbi.nlm.nih.gov/38603479/}{Spinal sensorimotor adaptation}~\cite{lavaud2024}
is a constrained example, not evidence of unrestricted skill transfer or
validation of our proposed pathway. Our comparison draws on the abstract and
relevant author-manuscript passages; equivalence to the final published edition
and the complete methods and statistics remain unverified here.

Thought2Action relates a declared intention to an independently measured action.
Action2World follows the action to its consequence in the task environment;
Thought2World names the proposed composition. Body2World drew attention to the
consequence outside the body; Body2Action focuses on body mechanics and a
specified task outcome, while Brain2Motion focuses on the movement pathway.
These working labels identify related scopes for a test. They do not establish
interchangeable or working controllers.

The engineering route connects neural activity to musculoskeletal simulation
and then to a programmable model, toward NeuroTranslation. The proposed
movement-focused learner is a large movement model. System 3 thinking names
the proposed human--AI supervisory partnership around it. The supervisory
model would compare task strategies and diagnose where neural or mechanical
predictions miss independent measurements, then propose training or calibration
comparisons for the signal mapping and virtual body. The task learner would
perform the declared prediction or control; the supervisor would help choose
what to test. Neither role replaces the biological spinal intermediary or a
physical interface. Final held-out outcomes must remain outside rewards and
model selection.
\par\noindent Shaw and Nave's SSRN preprint,
\emph{Thinking---Fast, Slow, and Artificial}, and its
\href{https://executiveeducation.wharton.upenn.edu/thought-leadership/wharton-at-work/2026/05/thinking-fast-slow-and-artificially/}{institutional discussion}
provide conceptual context for this proposed human--AI role.~\cite{shawNaveSystem3Retained}
The comparison is limited to that framing; the complete preprint and institutional
article have not been reviewed here. Neither establishes a third anatomical brain
system or validates this programme's proposed learner or supervisor. Expert footage may supply
movement and tactical examples, with rights and camera uncertainty accounted
for; it supplies no measured expert EEG, EMG, muscle force or stimulation plan.

Action changes the outcome. Its sensory consequences inform perception and
the next intention. Independently measured movement would anchor the action
part of that loop; it would not by itself decode a unique sensation or
perception. Known stimuli, measured feedback and appropriate recordings or
reports would have to test those further mappings. The return path must
therefore be measured rather than drawn into existence. The full
Interoceptomics aim includes internal-body sensing and connectivity; external
touch and proprioceptive tasks contribute evidence but do not by themselves
solve that wider problem. A prediction, assisted execution and retained skill
without assistance answer different questions. Each wider claim must earn
its own independent comparison.

\subsection{Context connects. Grounding checks}
Within a reach, context has a concrete job: relate brain activity to muscle
activity in this task and at this time. Attention is one proposed way to
represent those relationships. The model compares a query with keys and uses the resulting weights to combine their values. Masks keep unavailable observations and future
samples out. The calculation weights information; it does not establish that
attention understands all context.
% Source registry IDs: context-native-01a12f62, context-attention2017-v7, context-mult2019, context-policy4
\par\noindent \href{https://arxiv.org/html/1706.03762v7}{Vaswani et al. (2017)}~\cite{vaswani2017}; \href{https://aclanthology.org/P19-1656/}{Tsai et al. (2019)}~\cite{tsai2019}. These methods motivate the computation; validation of the proposed paired system still requires the independent observations below.

A pattern can fit inside the model and still describe the wrong movement.
Grounding therefore means testing the declared prediction against qualified
observations under the same conditions. Independently measured movement would
check the outcome, while muscle activity would add an observation that
constrains the brain model. Even a matching reach cannot identify a unique
private experience. No acquired paired data or trained model is established
for this proposed loop; those comparisons still have to be made.

\subsection{Muscle activity tests the brain model}
Record brain and named-muscle activity during the same reach. In the proposed
paired trial, EEG would sample scalp electrical activity and surface EMG would
sample activity associated with the named muscles. Paired training could use
the muscle observation to constrain their relationship. Then withhold the
muscle target and movement reference from the EEG-only predictor: can the
permitted brain observations predict a specified muscle-activity observable
or movement?

Pairing and calibration remain proposed. Shared cues and movement artefacts
could give that predictor a shortcut, so the comparison needs controls for
both. The muscle observations would cover only the named muscles, and EMG
supplies neither afferent nerve activity nor direct force. Sensory claims
still need known stimuli and measured feedback. Several neural states could
remain compatible with even a well-matched pair; the pair could narrow a
task-specific explanation without revealing the complete hidden brain.
Checking the movement itself requires the next independent observation.
% Source registry IDs: context-native-01a12f62, context-policy4, overhaul-brantley2018
\par\noindent Brantley and colleagues reported simultaneous scalp EEG, lower-limb EMG and locomotion observations from 17 body-worn inertial measurement units (IMUs) in ten able-bodied people. This supplies a paired gait precedent; it does not establish the proposed hand-reach apparatus, afferent recordings or optical three-dimensional motion capture. Selected-file rights, clocks, calibration and movement-artifact handling still need qualification. \href{https://www.nature.com/articles/sdata2018133}{Brantley et al., 2018}.

\subsection{Measured motion tests the virtual body}
Measure the same hand and arm from the aligned left and right tripod views.
The proposed motion-capture relay-composite would align the two point maps
in a common clock and coordinate frame, then join them into one composite
while retaining each point's origin. Before that joined map can check a
prediction, its calibration, alignment and capture precision would have to
be independently verified.

The movement map would constrain the external movement model: surface motion
and pose in the task's physical frame. To explain that motion mechanically,
the internal musculoskeletal model would describe specified bones, joints
and soft tissues. Organism, developmental stage and task would determine the
appropriate anatomical sources. X-ray projections could inform bone and
joint relationships; appropriate MRI references could inform specified soft
tissues. Registration, rights and source-specific limits would need independent
checks. No scan set is selected for this proposed comparison. Anatomy still
leaves activation, forces and neural function to test.

With a qualified composite, the musculoskeletal-simulation physics-engine's
joints, muscles, contacts and physical limits would face a measured comparison.
Fit the model on training data, then compare its predicted reach with
held-out measured movement. A force claim would require an additional force
measurement. The body's form limits possible motion; measured function would
test the mechanics and show where they need refinement.
% Source registry IDs: context-native-01a12f62, context-component-names-01a12f81, context-policy4, context-opensim2018, context-myosuite-docs, context-precision-01a12f77
\par\noindent \href{https://journals.plos.org/ploscompbiol/article?id=10.1371/journal.pcbi.1006223}{Seth et al. (2018)}~\cite{opensim2018}; \href{https://myosuite.readthedocs.io/en/latest/}{MyoSuite documentation}~\cite{myosuiteDocsUndated}. The simulation tools supply methods for that comparison; the proposed capture and body models still need its independent validation.

A lifelike hand makes that comparison visible. It cannot establish biological
fidelity. Capture alignment, precision, mechanics and biological fidelity have
not been demonstrated by this proposal, and captured geometry cannot establish
hidden neural state, internal forces or unique subjective experience. The
measured reach would check a movement prediction within those limits.

\subsection{Two pairs. One conditional learner}
The two comparisons would meet in the proposed large movement model: EEG
observations paired with muscle activity, and measured movement compared
with predicted mechanics. Large movement model names the proposed family;
deep neural network describes a possible implementation. Large refers to
the intended scope, rather than a demonstrated parameter count. Encoders
would receive only the observations permitted for the named test. Two
branches would summarize the relationships and meet in a shared task
representation; separate output parts would predict each direction's targets.
% Source registry IDs: context-native-01a12f62, context-policy4, context-attention2017-v7, context-mult2019
\par\noindent \href{https://arxiv.org/html/1706.03762v7}{Vaswani et al. (2017)}~\cite{vaswani2017}; \href{https://aclanthology.org/P19-1656/}{Tsai et al. (2019)}~\cite{tsai2019}. These computational methods offer precedents for a design, leaving the proposed paired learner and its performance to be tested.

Signal features belong to branches, and branches belong to the learner.
That hierarchy describes the computation, not two anatomical thinking systems.
A common clock, task and physical frame would connect the branches; independent
measurements would check their errors. Form constrains the prediction, and
checked function would guide refinement. Several hidden states could still
fit the same reach, so neither a shared representation nor conditional
predictions would establish an exact inverse nervous system. The architecture
has not been selected, implemented, trained or independently accepted.
To choose it, we would have to compare it with simpler baselines on
appropriately reserved observations.

\subsection{Learn the movement, then test what transfers}

A reach can provide a training target even when no EEG accompanies it. Earlier movement supplies the input; a later observation supplies the target. A separate proposed comparison would ask whether learning that relationship in one person helps forecast another person's grasp-and-lift movement. The WAY-EEG-GAL descriptor provides a relevant continuous-recording candidate, alongside its distinct windowed trials and derived lift summaries.\citep{luciw2014way} This movement-only question needs no AllLifts Run/trial join or invented cross-person neural pairing. The original condition-specific five-weight study remains a separate design.

The mapping has to explain what two bodies share. Homologous structures share evolutionary anatomical origin; mechanically analogous structures can serve comparable roles without that origin. A human digit aperture and a robot gripper opening could support a task-space comparison after separate mechanical qualification; this human-only proposal would not test that robotic extension. Joint degrees of freedom, axes and signs, proportions, muscle--tendon roles, loads, inertia, contact and time scales remain body-specific. Similar geometry cannot establish equivalent dynamics or neural coding. Retargeted motion remains derived; tracking, estimation, prediction and simulation keep their different evidential roles.

For the proposed two-person comparison, the representation contains index-, thumb- and object-sensor positions relative to the wrist, digit aperture and causal wrist velocity: thirteen features, with thirteen availability masks. Each participant needs qualified position units, a proper wrist-frame rotation, orientation convention and sensor-attachment calibration. A positive median wrist-to-index-sensor distance, expressed in qualified metres, supplies the characteristic length from assigned training support only; it is not a bone measurement. Multiplying normalized velocity by the same fixed 0.1-second reference interval makes all thirteen features dimensionless. Physical velocity errors invert that scaling with length divided by 0.1 seconds. Missing features receive no target loss; an unavailable observation is not a measured zero. Full finger joints, individual muscle forces and subject-specific inertia are not supplied by these sensor features.

The continuous target would be the same observed feature set 100 milliseconds beyond the latest actually available filtered input time. The proposal retains fifty frames at 100 Hz and a causal 51-tap filter before decimating qualified 500-Hz kinematics by five. Its 50-millisecond linear-phase delay must remain in the physical timing account; shifting information backward would change the prediction claim. Coefficients, passband, timestamp convention, gaps and boundary handling still require qualification. Only named kinematic columns may enter: the documented kinematics matrix also contains force-plate signals.

Identity and independence answer different questions. Trustworthy versioned participant-release attribution, an immutable selected archive entry and a qualified capture can identify the participant record and input object without requiring one particular internal ID-field spelling. They cannot establish biological identity or independent original acquisitions. Shared samples, trims, derivatives and overlapping recordings must form indivisible acquisition groups before filtering, windows, normalization or role assignment. Different names, checksums or series numbers do not prove separation. The complete public non-ST candidate roster and these relations still need establishing; all 294 rows of the prior AllLifts table remain excluded from training, evaluation and assigning trials or acquisition groups, and ST, competition storage and the reserved LINK recordings remain unavailable.

Qualified P1 groups would supply source pretraining; P2 would supply target training and development. With $G\geq2$ independently established P2 acquisition groups, the fixed identifier-hash ranking assigns the first $\lfloor3G/5\rfloor$ to training and the remainder to development. Canonical original-parent identifiers, their serialization and the roster must be fixed before ranking or examining outcomes. This is a nominal 60/40 split, not chronological ordering or a promise of condition balance. If unresolved ancestry merges P2 into one group, the comparison stops; splitting its windows cannot manufacture an independent development group. No group or segment has been assigned by the proposed identity method.

Three matched arms would test the transfer. B0 pretrains on correctly mapped P2 training movement. B1 and B2 use identical P1 windows and order: B1 sets the index-role input to the thumb sensor, the thumb-role input to the object sensor and the object-role input to the index sensor, while B2 retains the correct correspondence. Both share the correct source calibration and scale before that corruption, keep the wrist anchor, and use the same correct target adapter and ordered target batches. Parent-balanced sampling must record unique windows and repetitions. The false mapping can be relearned or can alter feature availability; matched update counts do not match source diversity. These are limits on attributing any difference to anatomical correspondence, not reasons to replace the control after seeing its result.

The proposed learner is one GRU layer of hidden width 32, with 26 inputs and thirteen outputs. Each arm and fixed seed---17, 31 and 73---receives 200 pretraining and 100 target updates, batch size 16, identical per-seed initialization and the fixed last checkpoint. Adam uses learning rate 0.001, betas 0.9 and 0.999, epsilon $10^{-8}$ and zero weight decay, without clipping, scheduling or early stopping. Training minimizes mean squared residuals over observed target features; a batch with no observed targets is rejected. The resulting 2,700 updates are planned, not performed. No external checkpoint is selected; numerical updates, filtering, precision and actual workload limits remain to be checked independently.

For each development parent and seed, the proposed score takes the square root of summed squared errors divided by the count of observed target features, then averages parent scores without weighting by parent duration. A zero denominator rejects the comparison. Report feature counts, missingness, physical-unit errors, uncertainty and persistence at the same information endpoint, including every worse parent or feature. B2 would have to lower normalized error against both B0 and B1 in every fixed seed, with no applicable original criterion failing, to justify another investigation. This remains a two-person development comparison, not independent final confirmation, population transfer or biological teaching.

Broader motion collections could supply additional movement examples only after their exact assets, rights, task coverage and original-recording overlap are qualified. In \href{https://www.addbiomechanics.org/download_data.html}{AddBiomechanics}, measured markers and ground reactions must remain distinct from fitted kinematics and estimated inverse-dynamics torques. \href{https://amass.is.tue.mpg.de/license.html}{AMASS} requires checking the selected asset's research-use, non-commercial and redistribution restrictions. Copies across collections stay with their original acquisition. LaBraM's WAY source-family exposure is flagged, while exact checkpoint/participant/trial exposure remains unknown; reserving a recording later cannot erase earlier pretraining exposure. No such checkpoint is selected for this comparison.

Before adding X-ray/CT or MRI, inspect the anatomical sources already used by the selected musculoskeletal model. A scan is useful when it constrains a stated gap or supplies an independent check beyond that model. Its person or species, scale, registration, rights and uncertainty still matter. Another person's scan remains an external anatomical reference; a generic model does not implicitly measure this participant's geometry, activation or forces. No scan set is selected here.

Actual source eligibility, acquisition groups, units, frames, clock and calibration remain unresolved for this continuous comparison, so no new fit or successful transfer is reported. A useful movement prior could later be tested with actual paired neural/body observations and the permitted input masks. Its transferred targets would still not become simultaneous measurements. The next result must identify where the mechanical forecast improves, where it worsens and which measurement or assumption accounts for the difference; neural equivalence, causal physiology, sensory delivery and retained skill would require further experiments.

\subsection{Pretrain. Pair. Verify. Refine}
Those comparisons start with the observations we may use. Data rights,
calibration and held-out subjects, sessions and tasks determine which can
enter the study; resolve them before processing data or choosing a model.
Open-source code does not grant blanket permission to use datasets,
checkpoints or clinical interfaces. Qualified data could first train
representations for each recording type. LaBraM is one EEG encoder candidate,
not a complete brain-body writer. Synchronized paired observations could
then train the declared mapping from permitted inputs.

Paired prediction would give later refinement a baseline to beat. Before
optional reinforcement learning, post-training or task-specific fine-tuning,
the physical predictions and reward need checking against independent
outcomes. Offline reinforcement learning needs qualified coverage of actions,
transitions and rewards: labels alone cannot tell us what an untried action
would do. Backpropagation can update differentiable neural components, but
it cannot automatically pass through every physics engine; universal
differentiability is not established. Simpler baselines and frozen final
tests remain through every phase.
% Source registry IDs: context-native-01a12f62, context-policy4, context-labram2405, context-cql2020-abstract, context-opensim2018, context-myosuite-docs
\par\noindent \href{https://arxiv.org/html/2405.18765v1}{Jiang, Zhao and Lu (2024)}~\cite{labram2024}; \href{https://arxiv.org/abs/2006.04779v3}{Kumar et al. (2020)}~\cite{cql2020}; \href{https://journals.plos.org/ploscompbiol/article?id=10.1371/journal.pcbi.1006223}{Seth et al. (2018)}~\cite{opensim2018}; \href{https://myosuite.readthedocs.io/en/latest/}{MyoSuite documentation}~\cite{myosuiteDocsUndated}. These sources offer candidate methods for the sequence; validation of the proposed paired system still requires the experiment below.

No data acquisition, training, fine-tuning, RL rollout or accepted checkpoint
is established for this proposed loop. The next experiment must test what
the prediction and reward represent against independent outcomes. A reward
that the computer can calculate supplies no biological proof on its own.

\subsection{Respond quickly. Improve carefully}
\label{b2-frozen-response}
The proposed response path would use a frozen accepted version to estimate
the current state, predict declared targets and prepare encoded outputs for
a separately qualified interface. Permitted causal time windows would limit
it to the allowed past and present observations. Estimating the state of this
moment would not change learned weights. A separate adaptation path would
inspect errors, calibrate or recalibrate, and propose a candidate update;
an independent comparison and a versioned saved state whose restoration reproduces the specified computation would precede any replacement. Examining errors would not
silently change the running model, and no unchecked live weight promotion
would occur.

The fast-and-slow comparison is a Kahneman-like analogy for this proposed
division of work. It identifies neither two biological or anatomical systems
nor inherent thinking speeds for EEG, spatial data or other sensor modalities.
Real-time performance would have to be measured across acquisition, filtering,
alignment, inference, output preparation, transduction or actuation, feedback
and rendering. Report delay tails, jitter, missed deadlines, uncertainty and
performance under specified concurrent workloads; component inference time alone cannot
establish end-to-end performance. No accepted model, running system that adapts
under these rules or real-time demonstration is established here. The loop has
no measured end-to-end latency or accuracy guarantee.
% Source registry IDs: context-native-01a12f62, context-policy4

\subsection{Write sensation. Enact action}
\label{b2-output-interfaces}
The prediction ends at the computer. A separate interface would have to turn
it into sensation or action. Input and output describe direction relative to
that computer; afferent and efferent describe physiological direction relative
to the central nervous system. A computer output could drive a sensory
interface supplying input toward the nervous system, while a computer input
could record activity related to the body's motor output. Both boundaries
must remain explicit.

The sensory route would use encoders and a separately qualified visual, haptic
or other sensory transducer to deliver specified feedback. Direct neural
writing remains a separate later research interface, with its own device,
measurement, safety and clinical evidence. Direct cortical writing could have
a sensory role without being a peripheral afferent signal. Predicting an
EEG-like pattern would not deliver sensation: the sensory content, interface,
authority and actual measured response would need their own tests, including
conditions, error and delay.

The action route would send separately qualified encoded controls through a
controller to a virtual body, a robotic actuator or a later qualified biological
interface. Action encoders, controllers and actuators have distinct jobs in
that route. A robot command is an action-role output, not an anatomical efferent
nerve. Simulated action, physical actuation, retained learning and biological
causal claims each require their corresponding evidence. EEG and EMG recorders
neither supply sensation nor actuate a body. These sensory and action
subinterfaces belong within the unchanged five programme pillars.

The retained references supply no assessed methods comparison that makes EEG or EMG recorders equivalent to an implanted sensory writer or action device. The missing clinical-methods source is identified as an unresolved provenance gap in Appendix~\ref{app:reproducibility-availability}. No method or capability is inferred from that missing support. Physical writers, transducers and actuators remain unqualified and unpriced, outside the recording-and-virtual-modelling cost scope.

This recording-and-virtual-modelling proposal cannot test physical sensory
delivery or biological actuation. No such physical interface, stimulation,
acquisition or actuation is established here, and no loop, training or
physical-writing result is claimed.

A physical study would need its own qualified interface and procedure. This
proposal supplies no stimulation dose, implantation, electrode-placement,
acquisition or actuation procedure. No such participant study, stimulation, biological actuation or qualified physical training interface has been implemented or tested. Participant, device, safety and clinical evidence would be separate requirements.
% Source registry IDs: context-write-native-01a12f66, context-native-01a12f62, context-policy4

\subsection{Close the loop with evidence}
Return to the cup and test each prediction against the observation it was
asked to explain. The brain-to-body prediction would face withheld observations;
the body-to-recorded-neural estimate would face its own recorded target.
A sensory or action interface would need a further test of what it delivered,
under which conditions, with what error and delay. Two models can agree and
still be wrong; a high reward can have the same problem. Neither agreement nor
a successful prediction would, by itself, identify a unique inverse brain or
establish someone's private experience.

The programme's distinction remains that approximation gets us towards
plausibility and virtualization gets us to precision. Independent measurements,
repeatability, causal checks and stated error bounds must earn that precision.
The original held-out and real-reference controls would test the prediction;
retained-learning controls would test whether a gain remains after useful
guidance ends. Precision is not an automatic result, and the proposal establishes
no perfect biology.

The next experiment must report where each declared prediction works and where
it fails against its independent measured target. Its answer would determine
which part of the proposed loop needs refinement.
% Source registry IDs: context-native-01a12f62, context-write-native-01a12f66, context-precision-01a12f77, context-policy4, context-spec15, context-glossary13, b2b-circular-checklist-original

\subsection{Technical representation and mask}
The information boundary must survive the calculation. A model given tomorrow's
sample or today's withheld target could make a persuasive prediction of an
answer it was never meant to receive. The representation and mask specify
which comparison the proposed attention computation would actually perform.

Learned query, key and value representations give the proposed attention
calculation a precise form. An admitted mask would restrict which observations
may contribute; biological grounding still requires the independent comparison.
\[\operatorname{Attention}(Q,K,V)=\operatorname{softmax}\!\left(QK^{\mathsf{T}}/\sqrt{d_k}+M\right)V.\]
Here Q, K and V are queries, keys and values; M expresses the admitted mask. The dimension $d_k$ is the common feature width of each query and key vector. The equation is methodological context from the supplied technical proposal. Modality, channel, task and time are part of the observation contract. No weights, EEG encoder or gradient route are selected or executed.

\subsection{Permitted observations and causal modality routes}
Declare B=EEG, E=EMG, Y=independently measured movement, C=task/stimulus/context, and P=the source-bound physics model. Training may relate qualified paired observations. Inference uses only its explicitly admitted subset:

\begin{itemize}
\item EEG-only: permitted past/present B and declared context predict a specified   EMG/movement target. Withheld E, Y and future samples cannot enter the predictor.
\item Body-only: permitted E and/or Y predict specified recorded neural features as   a conditional distribution with uncertainty. This is not a unique inverse brain.
\item Hybrid: additional modalities require a separately named route and evaluation;   a hybrid result cannot be reported as EEG-only.
\item Simulation-derived inputs stay simulated. They do not replace independent   observations or inherit the biological validity of a rendered body.
\end{itemize}

Separate training signals from evaluation references. Cycle consistency can make two models agree while both are wrong. Shared cues, movement artefacts, subject identity, session overlap and preprocessing leakage need the original controls. Do not expose final held-out targets to model, physics or reward selection.

\subsection{Representation, supervised learning and policy refinement}
\begin{enumerate}
\setcounter{enumi}{-1}
\item Establish source rights, subjects/tasks, measurement units, calibration,    uncertainty, clock alignment and frozen subject/session/task splits. Audit    pretraining overlap before preprocessing and model selection.
\item Pretrain modality representations on the qualified inventory. LaBraM is one    candidate EEG encoder: neural-tokenizer training, masked neural-code prediction    and downstream fine-tuning. It is not a complete brain-body writer. Code,    checkpoint, dataset and individual-use rights remain separate qualifications.
\item Learn paired supervised mappings with named targets and explicit modality    masks. Compare appropriate simpler, cue-only and artefact baselines. Keep raw,    processed and model-derived values distinct.
\item Qualify physical predictions and the reward verifier against independent    measurements. A successful task, smaller calibrated movement error and valid    physical states are prospective reward ingredients, not completed results.    Declare loss/reward weights on training/development evidence. A computable    objective or high reward does not confer biological ground truth.
\item Where admitted action/transition/reward coverage supports it, evaluate optional    RL and post-training. Offline trajectory labels alone do not establish the    consequences of actions never supported by the data. Conservative Q-Learning    is a methodological candidate, not a safety or performance guarantee here.
\item Fine-tune for the declared task and qualify both directions on frozen final    tests, uncertainty, baselines and every original threshold. Keep failures.
\item Evaluate a version with its weights fixed. Any later update must be a separate, versioned comparison with independently checked results and reproducible saved state.
\end{enumerate}

An illustrative supervised objective would combine declared errors for
specified EMG, movement and recorded-neural targets with applicable physical
constraints and optional consistency terms. Each term needs its units,
measurement uncertainty and target definition before a weight is chosen.
Consistency can check whether the outputs fit together; an independent
observation checks whether they fit the measured outcome.

Backpropagation changes differentiable components of the computer model;
biological feedback is a different process. A physics engine in the path does
not automatically give the learner a gradient. Qualify the actual
OpenSim/MyoSuite rollout and optimization route, or a separately validated
surrogate, rather than assuming gradients through every route. No paired four-source training,
RL rollout or task checkpoint for this proposed loop is reported here.

